
现在核对公式 (3.3.10)。
等式 $\mathfrak h_1=\mathfrak m$
由定义~\ref{def:7-3-18} 立即得到。

为验证 (3.3.10.1)、(3.3.10.2)，
考察链
\[
\bigl([0,1]^{|\vec a|}\times
\mathcal M^{\mathrm{main},\mathfrak d^+}_{k+1}
(\beta;\vec P^+)\bigr)
^{\mathfrak s^+_{\vec a,\mathfrak d^+,k,\beta,\vec P^+}}
\]
的边界。
设 $(\vec t^{\,i},x_i)$ 是该零集中的一个发散序列，
其中
$\vec t^{\,i}=(t_1^i,\ldots,t_\ell^i)$。
则出现以下三种情形之一：
\par\noindent\textup{情形 1}\quad
$t_j^i$ 趋于 $1$。
\par\noindent\textup{情形 2}\quad
$t_j^i$ 趋于 $0$。
\par\noindent\textup{情形 3}\quad
$x_i$ 趋于
$\partial\mathcal M^{\mathrm{main},\mathfrak d^+}_{k+1}
(\beta;\vec P^+)$。

由相容条件~\ref{cond:7-3-17}，
情形 2 给出
$-\mathfrak h_k\circ\mathfrak b_{k+1,i}$。
$\mathfrak b_{k+1,i}$ 的定义见第 3.3.1 小节。

与第 3 章定理 3.5.11 的
$A_\infty$ 公式证明类似，
把情形 3 与
\[
(-1)^n\partial\bigl(
\mathfrak h_{\ell+1,\beta}
(\vec P^{\,1}\mathbin{\overline\otimes}\cdots
\mathbin{\overline\otimes}\vec P^{\,\ell+1})
\bigr)
\]
合并，再对 $\beta\in\Pi(M;L)$ 求和，
可得
\[
\sum_{\ell=1}^{k+1}
\mathfrak h_\ell\circ
\widehat{\mathfrak h}_{k-\ell+2}.
\]
这里回忆
$\partial=(-1)^n\mathfrak m_{1,\beta_0}$。

最后考察三种情形中最需仔细处理的情形 1。
若 $\mathfrak P=0$，
由相容条件~\ref{cond:7-3-13}、
\ref{cond:7-3-14} 以及
引理~\ref{lem:7-3-2} 证明中的论证，
情形 1 的边界是退化链。
不过，还须适当修改 (7.3.19) 中
$\mathfrak h$ 的定义，
使这一边界项作为奇异链严格等于零。
稍后再回到这一点。

若 $\mathfrak P\ne0$，则
\[
\vec P^{\,1}\mathbin{\overline\otimes}\cdots
\mathbin{\overline\otimes}\vec P^{\,\ell+1}
=P\mathbin{\overline\otimes}1
\quad\text{或}\quad
1\mathbin{\overline\otimes}P.
\]
由相容条件~\ref{cond:7-3-14}，
情形 1 相应的边界分别给出
\[
-P=\mathfrak P_2(P\mathbin{\overline\otimes}1)
\quad\text{或}\quad
P=\mathfrak P_2(1\mathbin{\overline\otimes}P).
\]
于是完成恒等式 (3.3.10) 的证明。

再用命题 3.3.8 的 $A_{n,K}$ 版本，
便知 $PD([L])$ 是该 $A_{n,K}$ 代数的同伦幺元。
定理~\ref{thm:7-3-1} 证明余下从
$A_{n,K}$ 过渡到 $A_\infty$ 的部分，
与第 7.2.8 小节的论证平行，故从略。

最后说明完成情形 1 的讨论
所需的 $\mathfrak h$ 的修改。
为此，须修改模空间
\[
\bigl([0,1]^{|\vec a|}\times
\mathcal M^{\mathrm{main},\mathfrak d^+}_{k+1}
(\beta;\vec P^+)\bigr)
^{\mathfrak s^+_{\vec a,\mathfrak d^+,k,\beta,\vec P^+}}
\]
方法是仿照命题~\ref{prop:7-2-170}
的证明，塌缩 (7.3.7) 中映射的纤维。
要使塌缩后的模空间仍带有
与各种遗忘映射等相容的三角剖分，
就须研究仓西邻域的自然分层，
以及它在遗忘映射
(7.3.7.1)、(7.3.7.2) 下的行为。
先描述这一分层；
沿用引理~\ref{lem:7-3-8} 的记号。
